\documentclass[10pt]{article}

% One-page summary of the frozen-support LS and fixed-ridge comparison.
% Linear ridge uses identification-only column norms and lambda 1e-5.
% DOMP supports, PNNN training, and inference-operation counts stay fixed.

\usepackage[margin=1.20cm]{geometry}
\usepackage{booktabs,amsmath,array,tabularx}
\setlength{\parindent}{0pt}
\setlength{\parskip}{3pt}
\setlength{\emergencystretch}{2em}

\title{\textbf{Full-Signal Comparison of Independent PN-IQ, GMP, and Sparse PNNN}}
\author{}
\date{}

\begin{document}
\maketitle

\paragraph{Objective and protocol.}
Every linear model uses the same 19,661 identification rows (approximately
4\% of the capture), its previously frozen DOMP support and regressor matrix,
and all 491,520 rows for full-signal evaluation. Ridge uses the fixed value
$\lambda=10^{-5}$ without validation selection or cross-validation. Each
column norm is calculated on identification only; coefficients are then
returned to the original regressor scale. Full-signal evaluation includes the
identification samples and is not an independent generalization test.

\paragraph{Main results.}
NMSE is $10\log_{10}(\sum |y-\widehat y|^2/\sum |y|^2)$. The PNNN entry is the
unchanged frozen network result; Ridge is not applicable to that nonlinear
model.

\begin{center}
\scriptsize
\renewcommand{\arraystretch}{1.16}
\setlength{\tabcolsep}{3.2pt}
\begin{tabularx}{\linewidth}{@{}
  >{\hsize=.95\hsize\raggedright\arraybackslash}X
  r r r r
  >{\hsize=1.05\hsize\raggedright\arraybackslash}X@{}}
\toprule
Model & Parameters & \shortstack{NMSE LS\\(dB)} &
\shortstack{NMSE Ridge\\$10^{-5}$ (dB)} &
\shortstack{FLOPs/\\sample} & Additional operations/sample \\
\midrule
Independent PN-IQ & 358 & -38.534 & -36.940 & 1062 &
12 complex magnitudes, 2 divisions \\
Independent I/Q without PN & 394 & -37.810 & -36.417 & 1086 &
12 complex magnitudes \\
Parameter-matched GMP & 358 & -38.170 & -36.686 & 1899 &
14 complex magnitudes \\
Sparse PNNN N12 & 358 & -35.900 & -- & 876 &
14 complex magnitudes, 2 divisions, 12 sigmoid evaluations \\
\bottomrule
\end{tabularx}
\end{center}

\begin{center}
\begin{minipage}{0.97\linewidth}
\footnotesize
\textit{Ridge changes only coefficient estimation: it does not change support,
parameter count, or inference FLOPs. The sparse-PNNN FLOP count assumes that
zero weights are skipped; executing the original full matrices costs 2252
FLOPs/sample. A complex magnitude includes one square root.}
\end{minipage}
\end{center}

\paragraph{Complex-equivalent reference.}
Separately, the 100-term complex GMP baseline and its coupled PN-IQ real
representation remain numerically equivalent for both LS and fixed Ridge. The
full-signal Ridge relative prediction error is $7.989\times10^{-13}$. This
statement does not refer to the 179-term parameter-matched GMP in the table.

\paragraph{Quantitative conclusions.}
Under LS, Independent PN-IQ and parameter-matched GMP both use 358 real
parameters; PN-IQ improves NMSE by approximately 0.364 dB and requires 837
fewer FLOPs/sample. The no-PN control uses 394 parameters, 36 more than PN-IQ,
yet is approximately 0.723 dB worse under LS. This comparison indicates that
the observed gain is not explained by independent I/Q freedom alone.

With the fixed Ridge value, the three linear rows reach -36.940, -36.417, and
-36.686 dB, respectively. Ridge therefore does not improve this capture at
$10^{-5}$, but its regularized estimates preserve every frozen support and
inference cost. Sparse PNNN N12 has the lowest analytical count when zero
weights are skipped, while PN-IQ has the best NMSE in both displayed linear
fits. GMP and the no-PN control avoid the 12 sigmoid evaluations of the sparse
network. Analytical FLOPs do not establish latency, power, or hardware cost.

\end{document}
